\chapter{The BV Master Formalism for Gauge Effective Actions}
\label{ch:bv-master-formalism}
\ConventionsMixed


\section{Why Antifields Are Introduced}

The BRST chapter introduced sources \(K^{a\mu}\) and \(L^a\) for the nonlinear
variations \(sA_\mu^a\) and \(sc^a\).  Their role is to serve as coordinates
dual to the BRST variations, thereby making the Slavnov--Taylor identity a
functional equation for the 1PI effective action.

The Batalin--Vilkovisky formalism is the intrinsic version of this
construction.  It assigns a dual coordinate, called an antifield, to every
field in the gauge-fixed or ungauge-fixed complex.  The antibracket pairs a
field with its antifield, and the master equation is the single equation that
encodes gauge invariance, BRST nilpotency, gauge fixing, and the compatible
renormalization identities.

The framework in this chapter is perturbative and local.  We work with a
regulated field theory or, equivalently for algebraic purposes, with the
graded algebra of local functionals in finitely many jets of fields.  Infinite
dimensional functional derivatives are therefore shorthand for the
corresponding regulated or local-jet operation.  When an identity involves the
BV Laplacian, the regulator and measure are part of the data.

\section{Fields, Antifields, and Gradings}

Let \(\{\Phi^A\}\) denote the full set of fields in the gauge-theory complex.
The index \(A\) includes discrete labels, spacetime indices, Lie-algebra
indices, and the spacetime point in condensed notation.  Examples are
\[
  A_\mu^a,\qquad c^a,\qquad \bar c^a,\qquad B^a,
  \qquad \psi,\qquad \bar\psi .
\]
Each field has a Grassmann parity
\[
  \epsilon_A\in\mathbb Z/2\mathbb Z
\]
and a ghost number \(\operatorname{gh}(\Phi^A)\).  The antifield paired with
\(\Phi^A\) is denoted by \(\Phi^*_A\) and is assigned
\[
  \epsilon(\Phi^*_A)=\epsilon_A+1\pmod 2,
  \qquad
  \operatorname{gh}(\Phi^*_A)=-1-\operatorname{gh}(\Phi^A).
\]
Thus antifields have opposite parity and complementary ghost number.  For the
Yang--Mills nonminimal sector,
\[
\begin{array}{c|cccc}
  \Phi^A & A_\mu & c & \bar c & B\\ \hline
  \operatorname{gh}(\Phi^A) & 0 & 1 & -1 & 0\\
  \operatorname{gh}(\Phi^*_A) & -1 & -2 & 0 & -1
\end{array}
\]
and the parities are reversed in the antifield row.

Let \(\mathcal F_{\rm BV}\) be the graded-commutative algebra of local
functionals in \(\Phi^A,\Phi^*_A\) and finitely many of their spacetime
derivatives, modulo integration by parts when integrated local functionals are
being compared.  A functional \(F\in\mathcal F_{\rm BV}\) has parity
\(|F|\) and ghost number \(\operatorname{gh}(F)\) when it is homogeneous.
Left and right variational derivatives are defined by
\[
  \delta F
  =
  \int
  \delta\Phi^A\,
  \frac{\delta^L F}{\delta\Phi^A}
  =
  \int
  \frac{\delta^R F}{\delta\Phi^A}\,
  \delta\Phi^A ,
\]
and similarly for antifields.  The distinction is necessary because odd
variables do not commute with the variation symbols.

\section{The Antibracket}

\begin{definition}[BV antibracket]
For homogeneous functionals \(F,G\in\mathcal F_{\rm BV}\), the BV
antibracket is
\begin{equation}
  (F,G)
  =
  \int
  \left[
  \frac{\delta^R F}{\delta\Phi^A}
  \frac{\delta^L G}{\delta\Phi^*_A}
  -
  \frac{\delta^R F}{\delta\Phi^*_A}
  \frac{\delta^L G}{\delta\Phi^A}
  \right]\dd^D x ,
  \label{eq:bv-antibracket}
\end{equation}
with summation over all field labels \(A\).  The spacetime integral is absent
when the same formula is applied to a finite-dimensional regulator.
\end{definition}

The bracket has parity and ghost number
\[
  |(F,G)|=|F|+|G|+1,
  \qquad
  \operatorname{gh}(F,G)=\operatorname{gh}(F)+\operatorname{gh}(G)+1.
\]
It satisfies
\[
  (F,G)
  =
  -(-1)^{(|F|+1)(|G|+1)}(G,F)
\]
and the odd Jacobi identity
\[
  (F,(G,H))
  =
  ((F,G),H)
  +
  (-1)^{(|F|+1)(|G|+1)}(G,(F,H)).
\]
It is an odd biderivation.  Explicitly,
\[
  (F,GH)
  =
  (F,G)H
  +
  (-1)^{(|F|+1)|G|}G(F,H),
\]
and
\[
  (FG,H)
  =
  F(G,H)
  +
  (-1)^{|G|(|H|+1)}(F,H)G .
\]
The elementary brackets are
\[
  (\Phi^A(x),\Phi^*_B(y))=\delta^A{}_B\,\delta^{(D)}(x-y),
  \qquad
  (\Phi^A,\Phi^B)=0,
  \qquad
  (\Phi^*_A,\Phi^*_B)=0,
\]
with the delta distribution understood in the same regulated or distributional
sense as the fields themselves.

\section{The Classical Master Equation}

\begin{definition}[Classical BV action]
A classical BV action is an even local functional \(S_{\rm BV}\in
\mathcal F_{\rm BV}\) of ghost number zero satisfying the classical master
equation
\begin{equation}
  \frac12(S_{\rm BV},S_{\rm BV})=0 .
  \label{eq:classical-master-equation}
\end{equation}
It is required to reduce to the classical gauge-invariant action when all
ghosts and antifields are set to zero.
\end{definition}

The associated BV differential is
\[
  s_{\rm BV}F=(F,S_{\rm BV}).
\]

\begin{proposition}[Nilpotency from the master equation]
If \(S_{\rm BV}\) satisfies \eqref{eq:classical-master-equation}, then
\[
  s_{\rm BV}^2=0 .
\]
\end{proposition}

\begin{proof}
Using the odd Jacobi identity and the evenness of \(S_{\rm BV}\),
\[
  ((H,S_{\rm BV}),S_{\rm BV})
  =
  \frac12(H,(S_{\rm BV},S_{\rm BV}))
\]
for every functional \(H\).  The right hand side vanishes by
\eqref{eq:classical-master-equation}.  Hence
\[
  s_{\rm BV}^2H=((H,S_{\rm BV}),S_{\rm BV})=0 .
\]
\end{proof}

This proposition is the algebraic reason the master equation is stronger than
a compact rewriting of a Ward identity.  It constructs a differential on the
space of fields, ghosts, antifields, and local functionals.

\section{Minimal Yang--Mills BV Action}

For a gauge theory whose BRST transformations close off shell, the minimal BV
action is obtained by coupling antifields to the BRST variations.  With the
Hermitian Yang--Mills conventions of the preceding chapter,
\[
  sA_\mu^a=(D_\mu c)^a,\qquad
  sc^a=-\frac12 f^a{}_{bc}c^b c^c .
\]
The minimal BV action for the gauge and ghost sector is
\begin{equation}
  S_{\rm min}
  =
  S_{\rm YM}[A]
  +
  \int_M
  \left[
  A^{*\mu}_a(D_\mu c)^a
  -
  \frac12 c^*_a f^a{}_{bc}c^b c^c
  \right]\dd^D x .
  \label{eq:minimal-yang-mills-bv-action}
\end{equation}
Matter fields are included by adding their classical action and antifield
couplings
\[
  \int_M
  \left[
  \psi^*\,(\ii c^a t_R^a\psi)
  +
  (-\ii\bar\psi\,c^a t_R^a)\,\bar\psi^*
  \right]\dd^D x ,
\]
with the order of odd factors fixed by the left/right derivative convention.

The master equation for \eqref{eq:minimal-yang-mills-bv-action} is the
combined statement that \(S_{\rm YM}\) is gauge invariant and that the BRST
transformations are nilpotent.  Let
\[
  R_\mu^a(A,c):=(D_\mu c)^a,
  \qquad
  R_c^a(c):=-\frac12 f^a{}_{bc}c^b c^c,
\]
and write
\[
  S_{\rm min}=S_{\rm YM}+S_1,
  \qquad
  S_1=\int_M\left(A^{*\mu}_a R_\mu^a+c^*_aR_c^a\right)\dd^D x .
\]
Give each antifield antifield number \(1\) and each field antifield number
\(0\).  Since \(S_{\rm YM}\) is independent of antifields and \(S_1\) is
linear in antifields,
\[
  \frac12(S_{\rm min},S_{\rm min})
  =
  (S_{\rm YM},S_1)+\frac12(S_1,S_1),
\]
and the two displayed terms have antifield numbers \(0\) and \(1\),
respectively.  The antifield-number-zero component is
\[
  (S_{\rm YM},S_1)
  =
  \int_M
  \frac{\delta^R S_{\rm YM}}{\delta A_\mu^a}
  R_\mu^a\,\dd^D x
  =
  \int_M
  \frac{\delta S_{\rm YM}}{\delta A_\mu^a}
  (D_\mu c)^a\,\dd^D x .
\]
Gauge invariance of \(S_{\rm YM}\) means that this expression vanishes when
the odd ghost \(c\) is replaced by an arbitrary even infinitesimal gauge
parameter \(\zeta\).  Because the identity is algebraic in the parameter and
its derivatives, the same coefficient identity gives the vanishing for \(c\).

For the antifield-number-one component, the antibracket calculation gives
\[
  \frac12(S_1,S_1)
  =
  \int_M
  \left[
    A^{*\mu}_a\,sR_\mu^a
    +
    c^*_a\,sR_c^a
  \right]\dd^D x ,
\]
where \(s\) acts on functions of \(A,c\) by
\[
  sX
  =
  \int_M
  \left[
    R_\mu^a\frac{\delta_L X}{\delta A_\mu^a}
    +
    R_c^a\frac{\delta_L X}{\delta c^a}
  \right]\dd^D x .
\]
The two coefficients are explicitly
\[
\begin{aligned}
  sR_\mu^a
  &=
  D_\mu R_c^a
  +
  f^a{}_{bc}R_\mu^b c^c
\\
  &=
  -\frac12 f^a{}_{bc}D_\mu(c^b c^c)
  +
  f^a{}_{bc}(D_\mu c)^b c^c
  =
  0,
\end{aligned}
\]
where \(D_\mu\) is an even derivation and the two terms in
\(D_\mu(c^b c^c)\) are equal after antisymmetry of \(f^a{}_{bc}\) and the odd
commutation of \(c\) with \(D_\mu c\).  Similarly,
\[
\begin{aligned}
  sR_c^a
  &=
  -\frac12 f^a{}_{bc}
  \left(R_c^b c^c-c^b R_c^c\right)
\\
  &=
  -\frac1{12}
  \left(
    f^a{}_{be}f^e{}_{cd}
    +
    f^a{}_{ce}f^e{}_{db}
    +
    f^a{}_{de}f^e{}_{bc}
  \right)c^b c^c c^d
  =
  0
\end{aligned}
\]
by the Jacobi identity.  There are no higher antifield-number components for
this off-shell closed Yang--Mills algebra.  Thus
\(\frac12(S_{\rm min},S_{\rm min})=0\) is exactly the pair of statements
\(sS_{\rm YM}=0\) and \(s^2A=s^2c=0\).

If the gauge algebra closes only on shell, or if there are reducibility
relations among gauge transformations, the action contains additional
antifield terms.  The defining condition remains
\(\frac12(S_{\rm BV},S_{\rm BV})=0\); the higher antifield terms are not
extra notation but the terms required for the BV differential to square to
zero off shell.

\section{Nonminimal Sector and Gauge Fixing}

Gauge fixing is implemented by adding a contractible nonminimal pair and then
choosing a Lagrangian submanifold in the BV field-antifield space.  The
cohomological content is carried by the same BV differential.

For Yang--Mills theory the nonminimal pair is \((\bar c^a,B^a)\), with
\[
  s\bar c^a=B^a,\qquad sB^a=0.
\]
The nonminimal BV term is
\[
  S_{\rm nm}
  =
  \int_M \bar c^*_a B^a\,\dd^D x .
\]
The total classical BV action is
\[
  S_{\rm BV}=S_{\rm min}+S_{\rm nm}.
\]

\begin{definition}[Gauge-fixing fermion]
A gauge-fixing fermion is an odd functional \(\Psi[\Phi]\) of ghost number
\(-1\) depending only on the fields.  It determines the Lagrangian
submanifold
\begin{equation}
  \Phi^*_A=\frac{\delta^L\Psi}{\delta\Phi^A}.
  \label{eq:gauge-fixing-lagrangian-submanifold}
\end{equation}
The gauge-fixed action is the restriction
\[
  S_\Psi[\Phi]
  =
  S_{\rm BV}\left[\Phi,\Phi^*=\frac{\delta^L\Psi}{\delta\Phi}\right].
\]
\end{definition}

For covariant Yang--Mills gauges,
\[
  \Psi
  =
  \int_M
  \bar c^a\left(
  \partial^\mu A_\mu^a+\frac{g_{\rm YM}^2\xi}{2}B^a
  \right)\dd^D x .
\]
Substituting \eqref{eq:gauge-fixing-lagrangian-submanifold} into
\(S_{\rm BV}\), and integrating by parts with the same boundary convention as
in the BRST chapter, gives
\[
  S_\Psi
  =
  S_{\rm YM}
  +
  \int_M
  \left[
  B^a\partial^\mu A_\mu^a
  +
  \frac{g_{\rm YM}^2\xi}{2}B^aB^a
  -
  \bar c^a\partial^\mu D_\mu^{ab}(A)c^b
  \right]\dd^D x .
\]
Thus the Faddeev--Popov action is the restriction of the BV action to a
gauge-fixing Lagrangian submanifold.  Different gauge choices correspond to
different Lagrangian submanifolds.  They define equivalent physical
cohomology when the BV integral has no boundary contribution and the quantum
master equation is preserved.

\begin{figure}[H]
  \centering
  \resizebox{0.92\textwidth}{!}{%
  \begin{tikzpicture}[>=Stealth,thick,node distance=2.2cm]
    \node[draw,rounded corners=2pt,fill=blue!8,align=center,minimum width=3.4cm,minimum height=1.1cm] (fields)
      {fields and ghosts\\\(\Phi^A\)};
    \node[draw,rounded corners=2pt,fill=purple!8,align=center,right=of fields,minimum width=3.4cm,minimum height=1.1cm] (anti)
      {antifields\\\(\Phi^*_A\)};
    \node[draw,rounded corners=2pt,fill=green!8,align=center,right=of anti,minimum width=3.6cm,minimum height=1.1cm] (master)
      {master equation\\\((S_{\rm BV},S_{\rm BV})=0\)};
    \node[draw,rounded corners=2pt,fill=orange!10,align=center,below=1.55cm of anti,minimum width=4.2cm,minimum height=1.1cm] (gf)
      {gauge fixing\\\(\Phi^*=\delta\Psi/\delta\Phi\)};
    \node[draw,rounded corners=2pt,fill=red!8,align=center,right=of gf,minimum width=4.2cm,minimum height=1.1cm] (eff)
      {1PI/Wilsonian identities\\BV master equation};
    \draw[->] (fields)--(anti) node[midway,above] {dualize};
    \draw[->] (anti)--(master) node[midway,above] {antibracket};
    \draw[->] (master)--(eff);
    \draw[->] (anti)--(gf);
    \draw[->] (gf)--(eff);
  \end{tikzpicture}
  }
  \caption{BV replaces separate sources for nonlinear BRST variations by
  antifields for all variables.  Gauge fixing is a choice of Lagrangian
  submanifold, while effective actions are constrained by the same master
  equation.}
  \label{fig:bv-field-antifield-master-map}
\end{figure}

\section{The Quantum Master Equation}

The classical master equation controls the algebra of gauge symmetry.  The
quantum functional integral also depends on a measure.  In a finite-dimensional
odd symplectic regulator with coordinates \(\Phi^A,\Phi^*_A\), the BV
Laplacian is
\begin{equation}
  \Delta_{\rm BV}
  =
  \sum_A
  (-1)^{\epsilon_A}
  \frac{\partial^R}{\partial\Phi^A}
  \frac{\partial^R}{\partial\Phi^*_A}.
  \label{eq:bv-laplacian}
\end{equation}
In field theory this operator requires a regulator because the two functional
derivatives act at the same point.  A choice of regularization and measure is
therefore part of the definition of the quantum BV problem.

\begin{definition}[Quantum master equation]
For Lorentzian weight \(\exp(\ii S/\hbar)\), the quantum master equation is
\begin{equation}
  \frac12(S,S)-\ii\hbar\,\Delta_{\rm BV}S=0.
  \label{eq:quantum-master-equation}
\end{equation}
Equivalently,
\[
  \Delta_{\rm BV}\exp(\ii S/\hbar)=0.
\]
\end{definition}

The exponential form is conceptually direct: the integrand is closed under the
divergence operator associated with the chosen BV measure.  Gauge independence
is then a version of Stokes' theorem: changing the gauge-fixing Lagrangian
submanifold changes the integral by the integral of a BV-exact term, provided
no boundary contribution is present.

If the regulated theory fails to satisfy \eqref{eq:quantum-master-equation},
the failure has ghost number \(+1\).  To first order in \(\hbar\), write
\begin{equation}
  \mathcal M(S)
  :=
  \frac12(S,S)-\ii\hbar\,\Delta_{\rm BV}S
  =
  \hbar\,\mathcal A+O(\hbar^2).
  \label{eq:qme-anomaly-defect}
\end{equation}
The BV Jacobi identity, the second-order identity for \(\Delta_{\rm BV}\), and
\(\Delta_{\rm BV}^2=0\) give the quantum consistency identity
\begin{equation}
  (S,\mathcal M(S))-\ii\hbar\,\Delta_{\rm BV}\mathcal M(S)=0.
  \label{eq:qme-defect-consistency}
\end{equation}
Indeed, for even \(S\), \((S,\frac12(S,S))=0\) and
\(\Delta_{\rm BV}(S,S)=2(\Delta_{\rm BV}S,S)\), so the two order-\(\hbar\)
terms in \eqref{eq:qme-defect-consistency} cancel.  Substituting
\eqref{eq:qme-anomaly-defect} and \(S=S_0+O(\hbar)\) into
\eqref{eq:qme-defect-consistency} gives
\[
  0
  =
  \hbar\,(S_0,\mathcal A)+O(\hbar^2).
\]
Equivalently, with \(s_{\rm BV}:=(\cdot,S_0)\),
\[
  s_{\rm BV}\mathcal A=(\mathcal A,S_0)=0 .
\]
Changing the local counterterm by \(\hbar B\) changes the anomaly by
\begin{equation}
  \mathcal A\longmapsto \mathcal A+s_{\rm BV}B
  =
  \mathcal A+(B,S_0) .
  \label{eq:qme-anomaly-counterterm-shift}
\end{equation}
Thus the obstruction class lies in ghost-number-one local BV cohomology.  This
is the antifield version of the anomaly classification by
\(H^{1,D}(s\mid d)\).  In the chiral gauge example of
Chapter~\ref{ch:volume-ii-20}, the representative is the
descent cocycle whose bottom component gives the consistent anomaly
\eqref{eq:consistent-descent-anomaly-representative}; adding a local
counterterm changes the representative by an \(s_{\rm BV}\)-exact cocycle,
matching \eqref{eq:qme-anomaly-counterterm-shift}.

\section{The 1PI Master Equation}

The antifields are also the natural sources for nonlinear BRST variations in
the 1PI formalism.  With a gauge-fixing Lagrangian submanifold chosen for the
fields, define
\[
  Z[J,\Phi^*]
  =
  \int_{\mathcal L_\Psi}
  [D\Phi]\,
  \exp\left\{
  \frac{\ii}{\hbar}
  \left(S_{\rm BV}[\Phi,\Phi^*]+\int J_A\Phi^A\right)
  \right\}.
\]
The external antifields are not integrated; they record the response to BRST
variations.  Let \(W[J,\Phi^*]\) be the connected generator and let
\[
  \varphi^A=\frac{\delta W}{\delta J_A},
  \qquad
  \Gamma[\varphi,\Phi^*]
  =
  W[J,\Phi^*]-\int J_A\varphi^A
\]
be the 1PI effective action in the Lorentzian convention used earlier.

Assume that the regularized measure and action obey the quantum master
equation and that the Legendre transform is well defined on the field
directions under consideration.  The Ward identity obtained by applying
\(\Delta_{\rm BV}\) to the integrand gives
\[
  \int
  \frac{\delta^R\Gamma}{\delta\varphi^A}
  \frac{\delta^L\Gamma}{\delta\Phi^*_A}
  \dd^D x
  =
  0 .
\]
Equivalently,
\begin{equation}
  \frac12(\Gamma,\Gamma)=0,
  \label{eq:1pi-master-equation}
\end{equation}
where the bracket now pairs the mean field \(\varphi^A\) with the antifield
\(\Phi^*_A\).  In the Yang--Mills variables of the preceding chapter,
\eqref{eq:1pi-master-equation} is the Slavnov--Taylor identity
\eqref{eq:slavnov-taylor-identity}, with \(A^*\) and \(c^*\) replacing the
special sources \(K\) and \(L\), and with the nonminimal term producing
\(B\,\delta\Gamma/\delta\bar c\).

\section{Wilsonian BV Pushforward}

The Wilsonian effective action is obtained by integrating out degrees of
freedom above a scale.  BV gives the condition under which this operation is
compatible with gauge symmetry.

Consider first a finite-dimensional odd symplectic vector space split as
\[
  \mathcal E=\mathcal E_{\rm low}\oplus\mathcal E_{\rm high},
  \qquad
  \Delta_{\rm BV}=\Delta_{\rm low}+\Delta_{\rm high}.
\]
Let \(\mathcal L_{\rm high}\subset\mathcal E_{\rm high}\) be a Lagrangian
subspace used to integrate out the high variables.  Define the effective
low-energy integrand by
\begin{equation}
  \exp(\ii S_{\rm eff}/\hbar)
  =
  \int_{\mathcal L_{\rm high}}
  \exp(\ii S/\hbar).
  \label{eq:bv-wilsonian-pushforward}
\end{equation}

\begin{proposition}[BV pushforward]
Assume
\[
  \Delta_{\rm BV}\exp(\ii S/\hbar)=0
\]
and assume the integral over \(\mathcal L_{\rm high}\) has no boundary
contribution for \(\Delta_{\rm high}\).  Then
\[
  \Delta_{\rm low}\exp(\ii S_{\rm eff}/\hbar)=0 .
\]
\end{proposition}

\begin{proof}
Using \(\Delta_{\rm BV}=\Delta_{\rm low}+\Delta_{\rm high}\),
\[
\begin{aligned}
  \Delta_{\rm low}\exp(\ii S_{\rm eff}/\hbar)
  &=
  \int_{\mathcal L_{\rm high}}
  \Delta_{\rm low}\exp(\ii S/\hbar)
\\
  &=
  -\int_{\mathcal L_{\rm high}}
  \Delta_{\rm high}\exp(\ii S/\hbar)
  =
  0 .
\end{aligned}
\]
The last equality is the BV Stokes statement for the chosen Lagrangian
integration cycle.
\end{proof}

This proposition is the regulated mathematical core of the Wilsonian BV
identity.  In a field theory with cutoff \(\Lambda\), the Wilsonian action
\(S_\Lambda\) is a scale-dependent local or quasi-local representative of the
BV pushforward.  Its master equation has the form
\[
  \frac12(S_\Lambda,S_\Lambda)_\Lambda
  -
  \ii\hbar\,\Delta_\Lambda S_\Lambda
  =
  0,
\]
where the bracket and Laplacian include the cutoff covariance and the
remaining low-mode measure.  Exact RG flow is compatible with gauge symmetry
when it preserves this equation.  Changes of cutoff scheme are then BV
canonical transformations together with possible local counterterm changes,
with the gauge symmetry encoded by the transformed master action.

\section{Counterterms, Anomalies, and Physical Operators}

Let \(S_{\rm BV}\) be a classical solution of the master equation and let
\[
  s_{\rm BV}=(\cdot,S_{\rm BV}).
\]
A first-order local deformation
\[
  S_{\rm BV}\longmapsto S_{\rm BV}+\epsilon C
\]
preserves the classical master equation to order \(\epsilon\) precisely when
\[
  s_{\rm BV}C=0,
  \qquad
  \operatorname{gh}(C)=0.
\]
If \(C=s_{\rm BV}B\), then the deformation is generated by a BV canonical
transformation and does not change the physical theory.  Thus invariant
counterterms are classified by ghost-number-zero BV cohomology of local
functionals, modulo total derivatives and canonical redefinitions.

An anomaly candidate is a ghost-number-one functional \(\mathcal A\) satisfying
\[
  s_{\rm BV}\mathcal A=0.
\]
It is removable by a local counterterm exactly when
\[
  \mathcal A=s_{\rm BV}B .
\]
This gives the precise cohomological meaning of anomaly cancellation in the
perturbative BV framework: the obstruction class must vanish in the relevant
local BV cohomology.

Physical local operators are treated in the same language.  An insertion
\(\mathcal O\) is compatible with the gauge symmetry when
\[
  s_{\rm BV}\mathcal O=0,
  \qquad
  \operatorname{gh}(\mathcal O)=0,
\]
and two insertions differing by \(s_{\rm BV}\mathcal X\) define the same
cohomology class.  After gauge fixing, this statement reduces to the BRST
cohomology of local observables described in the preceding chapter.  Before
gauge fixing, it is expressed without choosing a particular gauge slice.

\section{Logical Output}

The output of this chapter is the following framework.

\begin{enumerate}
  \item A gauge theory is represented on an enlarged graded field space with
  fields, ghosts, nonminimal variables, and antifields.
  \item The antibracket is the odd symplectic pairing between fields and
  antifields.
  \item The classical master equation constructs a nilpotent differential.
  \item Gauge fixing is restriction to a Lagrangian submanifold determined by
  a gauge-fixing fermion.
  \item The quantum master equation is the regulated statement of
  measure-compatible gauge symmetry.
  \item The 1PI and Wilsonian effective actions in gauge theory must satisfy
  the corresponding BV master identities.
  \item Counterterms, anomalies, and physical operators are local BV
  cohomology classes at the appropriate ghost numbers.
\end{enumerate}

This is the framework used below whenever gauge-theory effective actions are
discussed beyond elementary Faddeev--Popov perturbation theory.
